\newpage
\section{System Overview}

\begin{block}
\begin{center}
    The source code can be found under the follwing link:\\
    {\tiSocialGithub}\, \url{https://github.com/s-grundner/LVA-EIC-KV}
\end{center}
\end{block}

\subsection{Pinout}

The project features four voices implemented as individual squarewaves oscillators. The output of each oscillator is routed to the outputpins \texttt{uo[0]}\dots\texttt{uo[3]}. Output \texttt{uo[7]} yields a PWM Signal, which encodes the number of actively running oscillators. This output may be used to control the gain of external circuitry when mixing the oscillator outputs together. The project has one input \texttt{ui[3]} to which the \acs{MIDI} signal has to be supplied to. This pin is also connected to the \acs{UART}-TX output of the RP2040 Microcrontroller, which is populated on the TinyTapeout test \acs{PCB}. The RP2040 may be used to emulate a \acs{MIDI} controller or to interface the instrument via \acs{USB}. \autoref{fig:pinout} summarises the in- and outputs of the ASIC.

\begin{figure}[h]
    \centering
    \includegraphics[width=\textwidth]{bsb-abstract.png}
    \caption{Input and output visualisation and description}\label{fig:pinout} 
\end{figure}

\subsection{Block Diagram}

An abstract overview of the project implementation can be seen in \autoref{fig:blockdiagram}. The block-titles outside of each block repesent the Verilog module name coherent with the code-repository.

\begin{figure}[h]
    \centering
    \includegraphics[width=\textwidth]{bsb-detailed.png}
    \caption{Abstract Blockdiagram of the Synthesizer (\texttt{synth} module)}\label{fig:blockdiagram}
\end{figure}

The signal flow starts at the input, where the \acs{MIDI} messages are received via \acs{UART}. The bits of the input stream are collected into bytes, which are then fed into a \acs{MIDI} decoder. The \acs{MIDI} decoder extracts the note- and channel number and if a on or off command is received. 

The note number indexes a lookup table which returns a counterbase for the corresponding frequency. Each oscillator in the stack is listening on its own \acs{MIDI}-channel. The note-on/note-off commands only affect the oscillator listening on this channel. Upon recieving a note on command, the counterbase value at the input of an oscillator is latched and a counter starts counting repeatedly in this time-base. The output of the oscillator is then toggled at each counter overflow, which generates the squarewave output.

\subsection{Layout}

To demonstrate the density of the layout, an image of the generated GDS is shown in \autoref{fig:layout}. 

\begin{figure}[htbp]
    \centering
    \includegraphics[width=0.9\textwidth]{layout.png}
    \caption{Layout of the Tile}\label{fig:layout}
\end{figure}

The routing statistics generated by the TinyTapeout pipeline are listed in \autoref{tab:routing}. Especially a space utilisation of \SI{74.275}{\percent} shows, that space-optimization was crucial. Per default, the maximum density is set to \SI{60}{\percent}, which had to be adjusted to fit more features. The documentation of the TinyTapeout project stated, that a density of up to \SI{80}{\percent} worked in previous shuttles without issues. The layout as it is shown in \autoref{fig:layout} is comprised of logic cells which are predefined by the \acs{PDK}. Which logic cells are generated from the Verilog code can be seen in \autoref{tab:cells}.

\begin{table}[htbp]
\caption{Routing statistics}\label{tab:routing}
\begin{tabularx}{\textwidth}{YY}
\toprule
Utilisation / \si{\percent} & Wire length / \si{\micro\meter} \\
\midrule
$74.275$ & $23342$ \\
\bottomrule
\end{tabularx}
\end{table}

% Table Generated from GDS Markdown Report with pandoc
\begin{table}[htbp]
\caption{Used Logic cells of the PDK and the number of occurences}\label{tab:cells}
\begin{tabularx}{\textwidth}{
  p{(\linewidth - 4\tabcolsep) * \real{0.2}}
  p{(\linewidth - 4\tabcolsep) * \real{0.7}}
  p{(\linewidth - 4\tabcolsep) * \real{0.1}}@{}
}
\toprule
Category & Cells & Count \\ \midrule
Fill &
\href{https://skywater-pdk.readthedocs.io/en/main/contents/libraries/sky130_fd_sc_hd/cells/decap/README.html}{decap}
\href{https://skywater-pdk.readthedocs.io/en/main/contents/libraries/sky130_fd_sc_hd/cells/fill/README.html}{fill}
& 733 \\
Combo Logic &
\href{https://skywater-pdk.readthedocs.io/en/main/contents/libraries/sky130_fd_sc_hd/cells/a22o/README.html}{a22o}
\href{https://skywater-pdk.readthedocs.io/en/main/contents/libraries/sky130_fd_sc_hd/cells/a41o/README.html}{a41o}
\href{https://skywater-pdk.readthedocs.io/en/main/contents/libraries/sky130_fd_sc_hd/cells/and2b/README.html}{and2b}
\href{https://skywater-pdk.readthedocs.io/en/main/contents/libraries/sky130_fd_sc_hd/cells/a2bb2o/README.html}{a2bb2o}
\href{https://skywater-pdk.readthedocs.io/en/main/contents/libraries/sky130_fd_sc_hd/cells/o22ai/README.html}{o22ai}
\href{https://skywater-pdk.readthedocs.io/en/main/contents/libraries/sky130_fd_sc_hd/cells/a31o/README.html}{a31o}
\href{https://skywater-pdk.readthedocs.io/en/main/contents/libraries/sky130_fd_sc_hd/cells/a22oi/README.html}{a22oi}
\href{https://skywater-pdk.readthedocs.io/en/main/contents/libraries/sky130_fd_sc_hd/cells/o21ai/README.html}{o21ai}
\href{https://skywater-pdk.readthedocs.io/en/main/contents/libraries/sky130_fd_sc_hd/cells/a221o/README.html}{a221o}
\href{https://skywater-pdk.readthedocs.io/en/main/contents/libraries/sky130_fd_sc_hd/cells/o22a/README.html}{o22a}
\href{https://skywater-pdk.readthedocs.io/en/main/contents/libraries/sky130_fd_sc_hd/cells/o2111ai/README.html}{o2111ai}
\href{https://skywater-pdk.readthedocs.io/en/main/contents/libraries/sky130_fd_sc_hd/cells/a211o/README.html}{a211o}
\href{https://skywater-pdk.readthedocs.io/en/main/contents/libraries/sky130_fd_sc_hd/cells/or4b/README.html}{or4b}
\href{https://skywater-pdk.readthedocs.io/en/main/contents/libraries/sky130_fd_sc_hd/cells/a21oi/README.html}{a21oi}
\href{https://skywater-pdk.readthedocs.io/en/main/contents/libraries/sky130_fd_sc_hd/cells/a21o/README.html}{a21o}
\href{https://skywater-pdk.readthedocs.io/en/main/contents/libraries/sky130_fd_sc_hd/cells/o211a/README.html}{o211a}
\href{https://skywater-pdk.readthedocs.io/en/main/contents/libraries/sky130_fd_sc_hd/cells/and3b/README.html}{and3b}
\href{https://skywater-pdk.readthedocs.io/en/main/contents/libraries/sky130_fd_sc_hd/cells/a311o/README.html}{a311o}
\href{https://skywater-pdk.readthedocs.io/en/main/contents/libraries/sky130_fd_sc_hd/cells/o21a/README.html}{o21a}
\href{https://skywater-pdk.readthedocs.io/en/main/contents/libraries/sky130_fd_sc_hd/cells/o2bb2a/README.html}{o2bb2a}
\href{https://skywater-pdk.readthedocs.io/en/main/contents/libraries/sky130_fd_sc_hd/cells/o221a/README.html}{o221a}
\href{https://skywater-pdk.readthedocs.io/en/main/contents/libraries/sky130_fd_sc_hd/cells/a2111o/README.html}{a2111o}
\href{https://skywater-pdk.readthedocs.io/en/main/contents/libraries/sky130_fd_sc_hd/cells/and4b/README.html}{and4b}
\href{https://skywater-pdk.readthedocs.io/en/main/contents/libraries/sky130_fd_sc_hd/cells/o211ai/README.html}{o211ai}
\href{https://skywater-pdk.readthedocs.io/en/main/contents/libraries/sky130_fd_sc_hd/cells/or4bb/README.html}{or4bb}
\href{https://skywater-pdk.readthedocs.io/en/main/contents/libraries/sky130_fd_sc_hd/cells/a211oi/README.html}{a211oi}
\href{https://skywater-pdk.readthedocs.io/en/main/contents/libraries/sky130_fd_sc_hd/cells/or3b/README.html}{or3b}
\href{https://skywater-pdk.readthedocs.io/en/main/contents/libraries/sky130_fd_sc_hd/cells/a31oi/README.html}{a31oi}
\href{https://skywater-pdk.readthedocs.io/en/main/contents/libraries/sky130_fd_sc_hd/cells/o41ai/README.html}{o41ai}
\href{https://skywater-pdk.readthedocs.io/en/main/contents/libraries/sky130_fd_sc_hd/cells/o21ba/README.html}{o21ba}
\href{https://skywater-pdk.readthedocs.io/en/main/contents/libraries/sky130_fd_sc_hd/cells/a21bo/README.html}{a21bo}
\href{https://skywater-pdk.readthedocs.io/en/main/contents/libraries/sky130_fd_sc_hd/cells/a32o/README.html}{a32o}
\href{https://skywater-pdk.readthedocs.io/en/main/contents/libraries/sky130_fd_sc_hd/cells/a221oi/README.html}{a221oi}
\href{https://skywater-pdk.readthedocs.io/en/main/contents/libraries/sky130_fd_sc_hd/cells/o2111a/README.html}{o2111a}
\href{https://skywater-pdk.readthedocs.io/en/main/contents/libraries/sky130_fd_sc_hd/cells/and4bb/README.html}{and4bb}
\href{https://skywater-pdk.readthedocs.io/en/main/contents/libraries/sky130_fd_sc_hd/cells/o32a/README.html}{o32a}
& 228 \\
Tap &
\href{https://skywater-pdk.readthedocs.io/en/main/contents/libraries/sky130_fd_sc_hd/cells/tapvpwrvgnd/README.html}{tapvpwrvgnd}
& 225 \\
Flip Flops &
\href{https://skywater-pdk.readthedocs.io/en/main/contents/libraries/sky130_fd_sc_hd/cells/dfrtp/README.html}{dfrtp}
\href{https://skywater-pdk.readthedocs.io/en/main/contents/libraries/sky130_fd_sc_hd/cells/dfstp/README.html}{dfstp}
& 191 \\
Multiplexer &
\href{https://skywater-pdk.readthedocs.io/en/main/contents/libraries/sky130_fd_sc_hd/cells/mux4/README.html}{mux4}
\href{https://skywater-pdk.readthedocs.io/en/main/contents/libraries/sky130_fd_sc_hd/cells/mux2/README.html}{mux2}
& 180 \\
Misc &
\href{https://skywater-pdk.readthedocs.io/en/main/contents/libraries/sky130_fd_sc_hd/cells/conb/README.html}{conb}
\href{https://skywater-pdk.readthedocs.io/en/main/contents/libraries/sky130_fd_sc_hd/cells/dlymetal6s2s/README.html}{dlymetal6s2s}
\href{https://skywater-pdk.readthedocs.io/en/main/contents/libraries/sky130_fd_sc_hd/cells/dlygate4sd3/README.html}{dlygate4sd3}
& 98 \\
AND &
\href{https://skywater-pdk.readthedocs.io/en/main/contents/libraries/sky130_fd_sc_hd/cells/and3/README.html}{and3}
\href{https://skywater-pdk.readthedocs.io/en/main/contents/libraries/sky130_fd_sc_hd/cells/and2/README.html}{and2}
\href{https://skywater-pdk.readthedocs.io/en/main/contents/libraries/sky130_fd_sc_hd/cells/and4/README.html}{and4}
\href{https://skywater-pdk.readthedocs.io/en/main/contents/libraries/sky130_fd_sc_hd/cells/a21boi/README.html}{a21boi}
& 95 \\
Buffer &
\href{https://skywater-pdk.readthedocs.io/en/main/contents/libraries/sky130_fd_sc_hd/cells/clkbuf/README.html}{clkbuf}
\href{https://skywater-pdk.readthedocs.io/en/main/contents/libraries/sky130_fd_sc_hd/cells/buf/README.html}{buf}
& 79 \\
NOR &
\href{https://skywater-pdk.readthedocs.io/en/main/contents/libraries/sky130_fd_sc_hd/cells/nor2/README.html}{nor2}
\href{https://skywater-pdk.readthedocs.io/en/main/contents/libraries/sky130_fd_sc_hd/cells/xnor2/README.html}{xnor2}
& 59 \\
Inverter &
\href{https://skywater-pdk.readthedocs.io/en/main/contents/libraries/sky130_fd_sc_hd/cells/inv/README.html}{inv}
& 49 \\
OR &
\href{https://skywater-pdk.readthedocs.io/en/main/contents/libraries/sky130_fd_sc_hd/cells/or2/README.html}{or2}
\href{https://skywater-pdk.readthedocs.io/en/main/contents/libraries/sky130_fd_sc_hd/cells/xor2/README.html}{xor2}
\href{https://skywater-pdk.readthedocs.io/en/main/contents/libraries/sky130_fd_sc_hd/cells/or3/README.html}{or3}
\href{https://skywater-pdk.readthedocs.io/en/main/contents/libraries/sky130_fd_sc_hd/cells/or4/README.html}{or4}
& 48 \\
NAND &
\href{https://skywater-pdk.readthedocs.io/en/main/contents/libraries/sky130_fd_sc_hd/cells/nand2/README.html}{nand2}
\href{https://skywater-pdk.readthedocs.io/en/main/contents/libraries/sky130_fd_sc_hd/cells/nand3/README.html}{nand3}
\href{https://skywater-pdk.readthedocs.io/en/main/contents/libraries/sky130_fd_sc_hd/cells/nand2b/README.html}{nand2b}
& 32 \\
Clock &
\href{https://skywater-pdk.readthedocs.io/en/main/contents/libraries/sky130_fd_sc_hd/cells/clkinv/README.html}{clkinv}
& 10 \\ \midrule
\multicolumn{3}{l}{1069 total cells (excluding fill and tap cells)} \\ \bottomrule
\end{tabularx}
\end{table}